% !TEX root = ../main.tex
%% Sujet type bac n° 1 — sujet inédit OptiTAB (valeurs vérifiées : verifications/sujets_bac.py)

\sujetBac{1}{{Un détartrant à l'acide lactique : acide faible, titrage, incertitude}/9,{Le lancer du poids : mouvement dans un champ uniforme}/6,{Un capteur d'humidité : circuit RC}/5}

\begin{exoBac}{9}{Un détartrant à l'acide lactique}
L'acide lactique est produit par la fermentation du lait et par les muscles lors d'un effort intense. Biodégradable, il
entre aussi dans la composition de détartrants dits « écologiques ». On étudie la molécule, puis on contrôle
l'étiquette d'un détartrant, avant de revenir à l'acide lactique formé dans l'organisme.
\begin{donnees}
Acide lactique : $M=\qty{90,1}{g.mol^{-1}}$. Couple acide lactique / ion lactate, noté $\ce{AH}/\ce{A-}$ : $\pKa=3{,}86$ à
\qty{25}{\degreeCelsius}. Zones de virage : hélianthine 3,1 -- 4,4 ; bleu de bromothymol 6,0 -- 7,6 ; phénolphtaléine 8,2 -- 10,0.
\end{donnees}

\partieExo{1}{La molécule d'acide lactique}
\begin{center}
\setchemfig{atom sep=2.1em}\chemfig{-[:30](-[:90]OH)-[:-30](=[:-90]O)-[:30]OH}\par\smallskip
{\legende{Formule topologique de l'acide lactique.}}
\end{center}
\begin{docu}{1}{Quelques bandes d'absorption en spectroscopie infrarouge}
{\small\begin{tabular}{@{}llll@{}}
Liaison & O--H (alcool) & O--H (acide carboxylique) & C$=$O (acide carboxylique) \\
Nombre d'onde (\unit{cm^{-1}}) & 3200 -- 3400 & 2500 -- 3200 & 1700 -- 1725 \\
Allure & forte, large & forte, très large & forte, fine \\
\end{tabular}}
\end{docu}
\begin{questions}
\item Recopier la formule topologique de l'acide lactique, entourer ses groupes caractéristiques et nommer les familles
fonctionnelles correspondantes.
\item Citer les bandes d'absorption attendues dans le spectre infrarouge de l'acide lactique.
\item Écrire l'équation de la réaction de l'acide lactique avec l'eau.
\end{questions}

\partieExo{2}{Contrôler l'étiquette d'un détartrant}
\begin{docu}{2}{L'étiquette et le protocole}
L'étiquette indique : « acide lactique 45~\% », c'est-à-dire 45~g d'acide lactique pour 100~g de détartrant, dont la masse
volumique vaut \qty{1,10}{g.mL^{-1}}. On prépare une solution S en diluant 100 fois le détartrant, puis on titre $V_A=\qty{20,0}{mL}$
de S par une solution d'hydroxyde de sodium, $\ce{Na+}+\ce{HO-}$, de concentration $c_B=\qty{0,100}{mol.L^{-1}}$, en suivant le pH.
Verrerie disponible : béchers de 50 et \qty{100}{mL} ; éprouvettes graduées de 10 et \qty{100}{mL} ; pipettes jaugées de 1,0 ; 10,0 et
\qty{20,0}{mL} ; fioles jaugées de 50,0 ; 100,0 et \qty{1000,0}{mL}.
\end{docu}
\schema{bac-s1-ph}
\begin{questions}[resume]
\item Montrer que la concentration en acide lactique annoncée par l'étiquette vaut environ \qty{5,5}{mol.L^{-1}}. En déduire
la valeur du volume équivalent attendue.
\item Choisir, en justifiant, la verrerie nécessaire pour préparer \qty{100,0}{mL} de solution S.
\item Écrire l'équation de la réaction support du titrage.
\item Déterminer le volume équivalent par une méthode graphique, que l'on explicitera.
\item Calculer le pH qu'aurait la solution S si l'acide lactique était un acide fort. Conclure à l'aide de la courbe.
\item Retrouver graphiquement la valeur du $\pKa$ du couple, en justifiant la méthode.
\item Quel indicateur coloré faudrait-il choisir pour réaliser ce titrage sans pH-mètre ? Justifier.
\item \initiative Le titre massique annoncé par l'étiquette est-il vérifié ? On admet que l'incertitude-type relative sur le
titre massique mesuré vaut 1,1~\%.
\end{questions}

\partieExo{3}{L'acide lactique dans le muscle}
Lors d'un effort intense, les cellules musculaires produisent de l'acide lactique ; le pH de leur cytoplasme reste
voisin de 7,4.
\begin{questions}[resume]
\item Représenter le diagramme de prédominance du couple $\ce{AH}/\ce{A-}$. Quelle forme prédomine dans la cellule ?
\item Calculer le rapport $\dfrac{[\ce{A-}]}{[\ce{AH}]}$ dans la cellule. Pourquoi les physiologistes parlent-ils plutôt de « lactate » ?
\end{questions}
\end{exoBac}

\begin{exoBac}{6}{Le lancer du poids}
Au lancer du poids, l'athlète propulse une boule de masse $m=\qty{7,26}{kg}$ et de diamètre \qty{12}{cm}. Lors d'un essai, la
boule quitte la main à la hauteur $h=\qty{2,10}{m}$ au-dessus du sol, avec une vitesse $v_0=\qty{13,5}{m.s^{-1}}$ faisant l'angle
$\alpha=\ang{40}$ avec l'horizontale. On étudie le mouvement du centre de la boule dans le référentiel terrestre, supposé
galiléen, muni du repère de la figure ; l'origine des dates est l'instant du lâcher.
\schema{bac-s1-poids}
\begin{donnees}
$g=\qty{9,81}{m.s^{-2}}$. Force de frottement de l'air : $F=\frac12\,\rho\,C\,S\,v^2$, avec $\rho=\qty{1,2}{kg.m^{-3}}$, $C=0{,}47$ et $S$ l'aire
du disque de même diamètre que la boule.
\end{donnees}
\begin{questions}
\item Calculer le poids de la boule et la valeur maximale de la force de frottement de l'air. Justifier que l'on néglige
les frottements.
\item À l'aide de la deuxième loi de Newton, déterminer les coordonnées du vecteur accélération.
\item Établir les équations horaires $x(t)$ et $y(t)$ du mouvement.
\item Montrer que la trajectoire a pour équation $y=-0{,}0459\,x^2+0{,}839\,x+2{,}10$, avec $x$ et $y$ en mètres.
\item Calculer la hauteur maximale atteinte par la boule.
\item \initiative Déterminer la longueur du jet, distance horizontale entre le pied de l'athlète, placé en $O$, et le point
d'impact.
\item Par un raisonnement énergétique, calculer la vitesse de la boule juste avant l'impact.
\item \capacite{Capacité numérique}L'entraîneur utilise le programme suivant.
\begin{codecote}[0.6]
\begin{pythoncol}
from math import sin, cos, sqrt, radians
g, v0, h = 9.81, 13.5, 2.10
meilleur, xmax = 0, 0
for a in range(20, 61):
    vx = v0 * cos(radians(a))
    vy = v0 * sin(radians(a))
    t = (vy + sqrt(vy**2 + 2*g*h)) / g
    if vx * t > xmax:
        meilleur, xmax = a, vx * t
print(meilleur, round(xmax, 2))
\end{pythoncol}
\cote
\begin{consolecol}
42 20.57
\end{consolecol}
\end{codecote}
Que représente la variable \code{t} ? Interpréter l'affichage. L'angle optimal serait de \ang{45} si la boule partait du sol :
expliquer la différence.
\end{questions}
\end{exoBac}

\begin{exoBac}{5}{Un capteur d'humidité}
Une station météorologique d'intérieur mesure l'humidité relative de l'air, notée HR, à l'aide d'un capteur capacitif.
\begin{docu}{1}{Le capteur}
Le capteur est un condensateur dont l'isolant, un polymère, absorbe la vapeur d'eau de l'air ambiant. Sa capacité, en
picofarads, vaut $C=161+0{,}34\times\mathrm{HR}$, où HR est exprimée en pourcentage (de 10 à 90~\%).
\end{docu}
\begin{docu}{2}{La mesure}
Un microcontrôleur charge le capteur, initialement déchargé, à travers une résistance $R=\qty{1,0}{M\ohm}$, sous la tension
$E=\qty{3,3}{V}$, et enregistre la tension $u_C$ à ses bornes.
\schema{bac-s1-rc}
\end{docu}
\begin{questions}
\item Établir l'équation différentielle vérifiée par $u_C$ pendant la charge.
\item Vérifier que $u_C(t)=E\left(1-\mathrm{e}^{-t/\tau}\right)$, avec $\tau=RC$, est solution de cette équation et respecte la condition
initiale.
\item Déterminer graphiquement le temps caractéristique $\tau$, puis la capacité du capteur.
\item En déduire l'humidité relative de la pièce.
\item \initiative Au-delà de 70~\% d'humidité, des moisissures risquent de se développer : le microcontrôleur déclenche
alors une ventilation. Quelle valeur de $\tau$ doit-il prendre comme seuil ?
\item Une mesure complète, charge et décharge, dure environ $10\,\tau$. Combien de mesures par seconde le microcontrôleur
peut-il réaliser au plus ?
\end{questions}
\end{exoBac}
